%=============================================================
%  第6章 結論
%=============================================================
\chapter{結論}
\label{chap:conclusion}

\section{研究成果の要約}
\label{sec:conclusion-summary}

本研究では，日本上場企業 \num{3558} 社・15 年間の財務データを用いて，
財務クラスタリングと年度間遷移の分析を行った。主要な成果を以下に要約する。

\begin{enumerate}
  \item \textbf{クラスタリング}:
    PCA と K-Means，およびハンガリアン法の組み合わせにより，
    日本上場企業を5つの財務クラスターに分類した
    （C0:中小規模標準型，C1:小規模財務脆弱型，C2:中規模財務優良型，
    C3:大企業型，C4:小規模財務健全型）。

  \item \textbf{遷移の実態}:
    年度間の自己遷移率は平均 80.1\% であり，企業の財務クラスターは
    年度をまたいで高い安定性を持つ。

  \item \textbf{遷移要因}:
    XGBoost + SHAP 分析により，遷移の主要なトリガーは
    営業利益の変化（Mean $|\mathrm{SHAP}|$ = 0.311）と
    長期借入金の変化（0.297）であることが示された。
    財務変数単独でも AUC = 0.745 の予測精度を達成する。

  \item \textbf{遷移先の決定要因}:
    遷移先クラスターの予測において AUC = 0.947 を達成した。
    C1（財務脆弱型）への遷移は営業利益の悪化，
    C3（大企業型）への遷移は総資産の拡大が主因であることが示された。

  \item \textbf{パス依存性}:
    「現在クラスター」変数が予測において重要な役割を果たすことが確認され，
    財務クラスターの移行には強いパス依存性（経路依存性）が存在することが示された。

  \item \textbf{ロバストネス}:
    PCA 推定期間の相違による AUC 差分は最大 0.018 に留まり，
    分析結果の頑健性が確認された。
\end{enumerate}

\section{学術的貢献}
\label{sec:conclusion-contribution}

本研究の学術的貢献は以下のとおりである。

\begin{enumerate}
  \item ハンガリアン法を用いた年度間ラベル整合型クラスタリング手法を提案した。

  \item 日本上場企業全体を対象とする大規模パネルデータを用いて，
    財務クラスター遷移の実態を記述した初めての研究である。

  \item XGBoost + SHAP を財務クラスター遷移分析に適用し，
    遷移の有無と遷移先の双方を予測・説明する統合的なフレームワークを構築した。
\end{enumerate}

\section{限界と課題}
\label{sec:conclusion-limitations}

本研究にはいくつかの限界がある。

\begin{enumerate}
  \item クラスター数 $K=5$ の選択には主観的な判断が含まれる。
    今後，Bayesian Information Criterion や Gap statistic などの
    客観的指標に基づく選択を検討する余地がある。

  \item 本研究は財務データのみを用いており，経営者の意思決定，
    産業環境，マクロ経済条件などの非財務的要因を考慮していない。

  \item \TODO{その他の限界を記述する}
\end{enumerate}

\section{今後の研究方向}
\label{sec:conclusion-future}

\begin{enumerate}
  \item \textbf{非上場企業への拡張}:
    財務データの入手可能性の問題はあるが，
    非上場企業を含めた分析は研究の外部妥当性を高める。

  \item \textbf{産業別分析}:
    産業ごとにクラスター構造が異なる可能性があり，
    産業分類を考慮した分析が有益である。

  \item \textbf{因果推論への拡張}:
    本研究は予測（予測精度）に焦点を当てており，
    因果効果の識別は行っていない。
    差分の差分法（DID）や操作変数法を用いた因果推論への拡張が考えられる。
\end{enumerate}
